%%%%%%%%%%  scan idx 209  |  folio 196  %%%%%%%%%%
tel que
\[
 H_i(\mathbf f^{n'},A)\longrightarrow H_i(\mathbf f^n,A)
\]
soit nul.
\end{enumerate}

\textbf{1.2.} Nous allons utiliser (1.1) pour prouver (0.1). Soient $k$ un corps, $B=k[[x,y]]$, $\mathbf f$ le système $(x,y)$, $M=\coprod_{n\ 0}B/(\mathbf f^n)$, $A$ la $B$-algèbre $D_B(M)$ (EGA $0_{\mathrm{IV}}$ 18.2.3). On a trivialement
\[
 H_2^A(\mathbf f^n,A)\simeq H_2^B(\mathbf f^n,A).\qquad\text{\footnotemark}
\]
\footnotetext{L'exposant A ou B indique relativement à quel anneau on calcule le complexe de Koszul.}
On a d'autre part
\[
 H_2^B(\mathbf f^n,A)\simeq H_2^B(\mathbf f^n,B)\oplus H_2^B(\mathbf f^n,M)
\]
mais $H_2^B(\mathbf f^n,B)=0$ puisque $\mathbf f$ est une suite régulière (EGA III 1.1.4 et $0_{\mathrm{IV}}$ 15.1.20), donc on a
\[
 H_2^B(\mathbf f^n,A)\simeq H_2^B(\mathbf f^n,M),
\]
On a enfin (EGA II 1.1.3.3)
\[
 H_2^B(\mathbf f^n,M)\simeq H_B^0(\mathbf f^n,M),
\]
$H_B^0(\mathbf f^n,M)$ s'identifiant canoniquement au sous-module ${}_{(\mathbf f^n)}M$ de $M$ annulé par $(\mathbf f^n)$. D'où finalement :
\begin{equation}\tag{*}
 H_2^A(\mathbf f^n,A\simeq {}_{(\mathbf f^n)}M.
\end{equation}
Comme le morphisme de transition $H_2^A(\mathbf f^{n'},A)\longrightarrow H_2^A(\mathbf f^n,A)$ est induit par la multiplication par $\mathbf f^{\,n'-n}$ (EGA III 1.1.6), on vérifie immédiatement, à l'aide de (*), que \emph{le pro-objet} $n\mapsto H_2^A(\mathbf f^n,A)$ \emph{n'est pas nul}. On déduit alors de (1.1) qu'il existe un $A$-module injectif $I$ tel que $H^2_Y(X,\widetilde I)\ne0$, donc tel que $\widetilde I$ ne soit pas flasque (SGA 2 I 2.12). Remarquons en passant que la projection $\Spec(A)\longrightarrow\Spec(B)$ est un homéomorphisme,
